\section{Cycle graphs}
\label{sec:cycles}

For graphs of maximum degree two, all normalized nondegenerate series
with nonzero quadratic coefficient give the same Hilbert series. The
linear case is separate, as are the degenerate cases of
Proposition~\ref{prop:degenerate-cycles}.

Let $C_n$ be the $n$-cycle, $n\ge3$, and let
$f=u+\alpha u^2+\cdots$, with $\alpha\ne0$.
By Proposition~\ref{prop:max-degree-two} we may take $\alpha=1/2$. With the cycle oriented
so that $X_i=\varphi_{i+1}-\varphi_i$, this gives
\begin{equation}\label{eq:cycle-normalization}
Y_i=\varphi_{i+1}-\varphi_i-\varphi_i\varphi_{i+1},
\qquad i\in\Z/n\Z,
\end{equation}
in $\Phi_n=\K[\varphi_0,\ldots,\varphi_{n-1}]/(\varphi_i^2)$.
Put $F_d=\langle Y_0,\ldots,Y_{n-1}\rangle^d$ and
$H_n=\hs_{C_n}^f$.

\subsection{The coefficients through degree three}
\label{sec:cycle-low-coefficients}

The low-degree calculation can be made independently of the more
involved closing argument for the top term. We prove completeness of the
relations of word degree at most three, rather than infer their ranks
from finite computations. In this subsection write $x_i=\varphi_i$,
$\ell_i=x_{i+1}-x_i$ and $q_i=x_ix_{i+1}$; thus
$Y_i=\ell_i-q_i$. All indices are taken modulo $n$.

\begin{theorem}[The first three coefficients]\label{thm:cycle-low-coefficients}
For every $n\ge7$,
\begin{equation}\label{eq:cycle-low-coefficients}
 H_n(t)=1+nt+\left(\binom{n+1}{2}-1\right)t^2
              +\left(\binom{n+2}{3}-3n\right)t^3+O(t^4).
\end{equation}
In particular,
$H_n(t)\equiv\operatorname{tr}M(t)^n-t^2\pmod{t^4}$ for every $n\ge7$,
where $M(t)$ is the transfer matrix of~\eqref{eq:propagation-matrix}.
\end{theorem}

Let $S=\K[z_0,\ldots,z_{n-1}]$, with homogeneous component $S_j$, and
let $\operatorname{ev}:S\to\Phi_n$ send $z_i$ to $Y_i$. This map is
not homogeneous for ordinary edge degree. The distinction between a
homogeneous polynomial in the $z_i$ and its possibly inhomogeneous image
will be used throughout the proof.

\begin{lemma}[The homogeneous cubic kernel]\label{lem:cycle-cubic-kernel}
For $n\ge7$, the kernel of $\operatorname{ev}|_{S_3}$ is the
$2n$-dimensional space
\begin{equation}\label{eq:cycle-cubic-kernel}
 \Span\{z_i^3,\ z_i^2z_{i+1}+z_iz_{i+1}^2:i\in\Z/n\Z\}.
\end{equation}
\end{lemma}
\begin{proof}
The displayed polynomials are independent in $S_3$. They vanish under
evaluation because
\begin{equation}\label{eq:cycle-cubic-products}
 Y_i^2=-2q_i,\qquad
 Y_iY_{i+1}=q_i-x_ix_{i+2}+q_{i+1}.
\end{equation}
We prove that they account for every cubic relation by comparing ordinary
edge degrees, starting with the highest one.

A product $Y_iY_jY_k$ with distinct, pairwise nonadjacent indices has
edge-degree-six term $-q_iq_jq_k$. These monomials are distinct: their
supports are proper subsets of the cycle, since $n\ge7$, and each run
of consecutive vertices in such a support has a unique covering by
disjoint adjacent pairs. Every other cubic word contains a square or
an adjacent pair, and has edge degree at most four
by~\eqref{eq:cycle-cubic-products}. Thus the coefficients of all the
pairwise nonadjacent triples vanish in a relation.

Modulo~\eqref{eq:cycle-cubic-kernel}, the remaining words are represented
by adjacent doubles $Y_i^2Y_{i+1}$ and by
\[
 D_{ij}=Y_i^2Y_j\quad(i,j\text{ nonadjacent}),\qquad
 A_i=Y_iY_{i+1}Y_{i+2},\qquad
 B_{ij}=Y_iY_{i+1}Y_j,
\]
where $j\notin\{i-1,i,i+1,i+2\}$ in the last expression. Thus $B_{ij}$
has exactly one adjacent pair of generator indices. Its edge-degree-four
part is
\[
 -q_iq_j+x_ix_{i+2}q_j-q_{i+1}q_j.
\]
The middle monomial determines $(i,j)$ uniquely and does not occur in
any of the other types of words. To see this, its four-element support
has either one adjacent pair, or a run of three and an isolated vertex.
In the first case the pair $\{j,j+1\}$ is determined. In the second,
two descriptions would require the isolated vertex to be at distance
two from both endpoints of the run of three, which is possible only on
a cycle of length six. Neither type of support is a union of two
disjoint adjacent pairs. Comparing these monomials eliminates all
$B_{ij}$.

The edge-degree-four terms of $D_{ij}$ and $D_{ji}$ are both $2q_iq_j$;
that of $A_i$ is $-q_iq_{i+2}$. Different unordered nonadjacent pairs
give different supports, again by uniqueness of a covering by adjacent
pairs. Therefore all remaining cancellations in degree four are
spanned by $D_{ij}-D_{ji}$ and
$A_i+\tfrac12D_{i,i+2}$, together with the adjacent doubles, whose
images already have degree three. Direct multiplication identifies
their images, up to nonzero scalar multiples, with
\begin{equation}\label{eq:cycle-residual-cubics}
 \tau_i=x_ix_{i+1}x_{i+2},\qquad
 E_{ij}=q_i\ell_j-q_j\ell_i,\qquad
 Z_i=q_{i+2}\ell_i.
\end{equation}
Here one $E_{ij}$ is used for each unordered nonadjacent pair. Indeed,
$D_{ij}-D_{ji}=-2E_{ij}$ and
$A_i+\tfrac12D_{i,i+2}=Z_i$. It remains to prove the independence of
these ordinary homogeneous cubics.

Encode the coefficients of the $E_{ij}$ by numbers $c_{ij}=-c_{ji}$,
with $c_{ij}=0$ when $i,j$ are equal or adjacent. Let $b_i$ be the
coefficient of $Z_i$ in a proposed relation. For fixed $i$ and
$k\notin\{i-1,i,i+1,i+2\}$, the monomial $x_ix_{i+1}x_k$ has exactly
one adjacent pair. Its coefficient gives
\begin{equation}\label{eq:cycle-residual-coefficients}
 c_{i,k-1}-c_{i,k}-b_{i-2}\mathbf1_{\{k=i-2\}}=0.
\end{equation}
Taking successively $k=i+3,\ldots,i+n-3$ yields
\[
 c_{i,i+2}=c_{i,i+3}=\cdots=c_{i,i+n-3}=:u_i.
\]
If the cyclic distance between $i$ and $j$ is at least three, then
$u_i=-u_j$ by antisymmetry. The graph on $\Z/n\Z$ joining all such
pairs is connected and has an odd cycle. For $n=7$ its distance-three
edges form a seven-cycle. For $n=8$ they form a spanning cycle and the
vertices $0,3,6,1,4,0$ form an odd cycle. For $n\ge9$, steps of lengths
three and four give connectivity, and $0,3,6,0$ is a triangle. It follows
that every $u_i=0$. The distance-two coefficients vanish as well, since
$c_{i,i+2}=u_i$ and $c_{i,i-2}=-u_{i-2}$. Thus all $c_{ij}=0$.
Taking $k=i-2$ in~\eqref{eq:cycle-residual-coefficients} gives
$b_{i-2}=0$. Finally, the $\tau_i$ are distinct monomials. This proves
the independence in~\eqref{eq:cycle-residual-cubics} and the lemma.
\end{proof}

\begin{proof}[Proof of Theorem~\ref{thm:cycle-low-coefficients}]
Put $\ell=\sum_i z_i$, $q=\sum_i z_i^2$ and
$R_0=\ell-\tfrac12q$. Since $\sum_i\ell_i=0$ and $Y_i^2=-2q_i$,
we have $R_0(Y)=0$.

The $\ell_i$ span the $(n-1)$-dimensional space $W$ of linear forms
in the $x_i$ whose coefficient sum is zero. The ordinary multiplication
map $\operatorname{Sym}^2W\to\Phi_n^{(2)}$ is an isomorphism: the
squares $(x_i-x_j)^2=-2x_ix_j$ span its target and both spaces have
dimension $\binom n2$. Hence the kernel of the homogeneous evaluation
$S_2\to\Phi_n^{(2)}$, $z_i\mapsto\ell_i$, is $\ell S_1$.

Consider a relation of word degree at most three. Its constant term is
zero, and its linear part is a multiple of $\ell$. Subtract the
corresponding multiple of $R_0$. Its remaining quadratic part belongs
to $\ell S_1$, so subtract a linear multiple of $R_0$ to remove it.
The remaining relation is a homogeneous cubic and is described by
Lemma~\ref{lem:cycle-cubic-kernel}. The three resulting spaces are
independent, by their lowest polynomial degrees. Thus
\[
 \ker(\operatorname{ev}|_{S_{\le3}})
 =\K R_0\oplus S_1R_0\oplus\ker(\operatorname{ev}|_{S_3}),
 \qquad
 \dim F_3=\binom{n+3}{3}-3n-1.
\]

For completeness, $R_0$ is the only relation of word degree at most
two. After subtracting its linear part as above, any further relation
would be $\ell A$ with $A\in S_1$. Since $\ell(Y)=q(Y)/2$, the cubic
$qA$ would belong to~\eqref{eq:cycle-cubic-kernel}. For every $j$, choose
an index $i$ nonadjacent to $j$. The coefficient of $z_i^2z_j$ in $qA$
is the coefficient of $z_j$ in $A$, whereas every element
of~\eqref{eq:cycle-cubic-kernel} has coefficient zero there. Hence
$A=0$ and $\dim F_2=\binom{n+2}{2}-1$. There is no linear relation
among the $Y_i$: the only possible one is their sum, which equals
$-\sum_iq_i\ne0$. Therefore $\dim F_1=n+1$.
Subtracting successive cumulative dimensions proves
\eqref{eq:cycle-low-coefficients}.

Finally put $\mathrm{tr}_n=\operatorname{tr}M^n$. The matrix in
\eqref{eq:propagation-matrix} gives
$\mathrm{tr}_3=1+3t$, $\mathrm{tr}_5=1+5t+15t^2$ and
$\mathrm{tr}_7=1+7t+28t^2+63t^3$. The recurrence $L(\mathrm{tr}_n)=0$ implies
$[t]\mathrm{tr}_n=n$ for $n\ge3$ and $[t^2]\mathrm{tr}_n=\binom{n+1}{2}$ for $n\ge5$.
For $n\ge8$, the first difference of $[t^3]\mathrm{tr}_n$ is
$3\binom{n-1}{2}-2\binom{n-2}{2}$, the first difference of
$\binom{n+2}{3}-3n$. The value at $n=7$ proves the trace comparison.
\end{proof}

This proves the trace comparison through degree three, uniformly in
$n\ge7$. Conjecture~\ref{conj:cycles} in full requires every degree $d<n/2$.

\subsection{The top term}
\label{sec:cycle-top}

\begin{theorem}[The top term for cycles]\label{thm:cycle-top}
For every $n\ge3$,
\[
\deg H_n=\left\lceil\frac n2\right\rceil+1.
\]
More precisely, the highest-degree terms are
\begin{equation}\label{eq:cycle-top-terms}
\operatorname{lt}_tH_{2m+1}=t^{m+2}\quad(m\ge1),\qquad
\operatorname{lt}_tH_{2m}=(2^{m-1}+1)t^{m+1}\quad(m\ge2).
\end{equation}
\end{theorem}

The proof uses coefficient extraction to cut the cycle. The path
boundary calculation then describes the annihilators of its penultimate
filtration spaces. We give these annihilators explicitly.

\subsection{Dual boundary filtrations on paths}
\label{sec:dual-paths}

In $A_N=\K[x_0,\ldots,x_{N-1}]/(x_i^2)$ use the nondegenerate pairing
\begin{equation}\label{eq:coefficient-pairing}
\langle g,h\rangle_N=[x_0\cdots x_{N-1}]gh.
\end{equation}
For a subspace $U\subseteq A_N$, write $U^\perp$ for its annihilator.
For each $i$, put $A_{N-1}^{(i)}=\K[x_j:j\ne i]/(x_j^2)$ and
write $\partial_{x_i}:A_N\to A_{N-1}^{(i)}$ for coefficient extraction.
It satisfies
\[
 \partial_{x_i}(ab)=\partial_{x_i}(a)\,\bar b+
                  \bar a\,\partial_{x_i}(b),\qquad \bar a=a|_{x_i=0}.
\]
As in Section~\ref{sec:exact}, this is a derivation relative to
evaluation at $x_i=0$. When the target is included in $A_N$, it is not
an ordinary $A_N$-valued derivation. Set
\[
y_i=x_{i+1}-x_i-x_ix_{i+1},\qquad
B_N^d=\langle x_0,y_0,\ldots,y_{N-2},x_{N-1}\rangle^d,\qquad
P_N^d=\langle y_0,\ldots,y_{N-2}\rangle^d.
\]
For $N=1$, these mean $B_1^d=\langle x_0\rangle^d$ and $P_1^d=\K$.
The first filtration is closed at both endpoints; the second is open
at both. Let
\[
\Omega_{N-1}=\sum_{i=0}^{N-1}\prod_{j\ne i}x_j.
\]
The unfiltered algebra generated by the $y_i$ is
\begin{equation}\label{eq:open-path-algebra}
C_N=\K\oplus
\left\{\sum c_ix_i:\sum c_i=0\right\}
\oplus\bigoplus_{j\ge2}A_N^{(j)},
\qquad C_N^\perp=\K\Omega_{N-1}.
\end{equation}
Indeed, replacing the $y_i$ by the invertible series giving
$x_{i+1}-x_i$ does not change their algebra. These differences span all
linear forms of coefficient sum zero. Their squares yield every
$x_ix_j$, $i\ne j$, and multiplication by further differences gives all
higher square-free monomials. This also proves the last assertion.
The same description, with $x_i=\varphi_i$, applies to $\CC_{C_N}$.

For $N=2r-1$, define
\begin{equation}\label{eq:dual-path-V}
\mathcal V_r=\Span\left\{
 \prod_{j=0}^{r-2}u_j:
 u_j\in\{1,x_{2j+1},x_{2j}-x_{2j+2}-x_{2j}x_{2j+2}\}
\right\};
\end{equation}
empty products give $\mathcal V_1=\K$.

\begin{lemma}[Dual path calculation]\label{lem:dual-path}
For $r\ge1$,
\begin{equation}\label{eq:dual-closed-path}
(B_{2r-1}^{r-1})^\perp=\mathcal V_r,
\qquad \dim\mathcal V_r=3^{r-1}.
\end{equation}
For $r\ge2$,
\begin{equation}\label{eq:dual-odd-open-path}
(P_{2r-1}^{r})^\perp
 =\K\Omega_{2r-2}\oplus
 (1-x_1)(1+x_{2r-3})\mathcal V_{r-1},
\end{equation}
where the second summand uses the interior variables
$x_1,\ldots,x_{2r-3}$. Moreover, $P_{2r-1}^{r+1}=C_{2r-1}$
and $P_{2r}^{r+1}=C_{2r}$; in particular
\begin{equation}\label{eq:dual-even-open-path}
(P_{2r}^{r+1})^\perp=\K\Omega_{2r-1}.
\end{equation}
\end{lemma}

\begin{proof}
First compute dimensions from the transfer matrix
$M$ in~\eqref{eq:propagation-matrix}. Besides the closed evaluation
$\mathsf c$, use
\[
\tau(e_0)=2,\qquad \tau(e_1)=1+t,\qquad \tau(e_2)=1+t^2.
\]
Then
\begin{equation}\label{eq:dual-path-series}
H(B_N)=\mathsf c(M^{N-1}e_1),\qquad
H(P_N)=1+t\sum_{j=0}^{N-2}\tau(M^je_1).
\end{equation}
For the second formula, use $[0,\infty]$ as an auxiliary symbol:
$1$ is present, but its boundary multiple never occurs in the open
filtration. This is not an element of $\mathcal S$, whose states have
two finite births. For this calculation only, extend the shift convention
by $t^a[p,\infty]=[p+a,\infty]$.
The same two-variable calculation as in
Lemma~\ref{lem:boundary-transition} gives
$T[0,\infty]=[0,\infty]+te_1$.
Iterating and evaluating an open state by $t^p+t^q$, with
$t^\infty=0$, proves the formula.
Write $M^{2j}e_1=A_je_0+B_je_1+C_je_2$. From
\[
M^2e_0=2te_0+2e_2,\qquad M^2e_1=3te_1+e_2,\qquad
M^2e_2=t^2e_0+te_1+(1+t)e_2
\]
we get $(A_0,B_0,C_0)=(0,1,0)$ and
\[
A_{j+1}=2tA_j+t^2C_j,\qquad
B_{j+1}=3tB_j+tC_j,\qquad
C_{j+1}=2A_j+B_j+(1+t)C_j.
\]
Induction on $j$ therefore gives
\[
\deg A_j\le j,\quad \deg C_j\le j-1,\quad
B_j=3^jt^j+\text{terms of lower degree};
\]
in the inductive step, $tC_j$ has degree at most $j$, so it does not
affect the leading term of $B_{j+1}$.
Since $H(B_{2r-1})=\mathsf c(M^{2r-2}e_1)$ by~\eqref{eq:dual-path-series},
and $H(P_N)$ involves in addition the odd powers $M^{2j+1}e_1$, this gives
\[
\deg H(B_{2r-1})=r,\quad [t^r]H(B_{2r-1})=3^{r-1},
\]
and, for $r\ge2$,
\[
\deg H(P_{2r-1})=r+1,\quad
[t^{r+1}]H(P_{2r-1})=3^{r-2},\qquad
\deg H(P_{2r})\le r+1.
\]
Together with~\eqref{eq:open-path-algebra}, these give the needed
annihilator dimensions and the exhaustion assertions.

To identify the subspace in~\eqref{eq:dual-closed-path}, append two
variables $b,c$ to a path whose last variable is $a$.
The three new factors in~\eqref{eq:dual-path-V} are
$1,b,a-c-ac$. For a single terminal generator one has
\begin{equation}\label{eq:path-extraction-table}
\left([bc]u\ell\right)_{
 u=1,b,a-c-ac;\ \ell=y(a,b),y(b,c),c}
 =\begin{pmatrix}0&-1&0\\0&1&1\\-1&1&0\end{pmatrix}.
\end{equation}
With no terminal generator, extraction is zero; with one, it is a
constant, by~\eqref{eq:path-extraction-table}. With two or more, it
involves only the variable $a$, so it lies in $\Span\{1,a\}$, and $a$ is
the closing generator of the shorter path. In each case extraction lowers
the closed-path filtration bound by at least one. Starting with
$(B_1^0)^\perp=\K$, induction proves that every element of $\mathcal V_r$
annihilates $B_{2r-1}^{r-1}$.
The three copies of $\mathcal V_{r-1}$ are independent: in
$h_0+bh_1+(a-c-ac)h_2$, the coefficient of $c$ is $-(1+a)h_2$;
the factor $1+a$ is a unit. Then the coefficient of $b$ gives $h_1=0$,
and finally $h_0=0$. Thus $\dim\mathcal V_r=3^{r-1}$, which is the
full annihilator dimension.

For~\eqref{eq:dual-odd-open-path}, let
$g=(1-x_1)(1+x_{2r-3})h$, with $h\in\mathcal V_{r-1}$.
A product contributing to the pairing must use both outer generators.
Extracting $x_0$ from $(1-x_1)y_0$ and $(1-x_1)y_0^2$ gives
$-1$ and $-2x_1$, respectively; extracting $x_{2r-2}$ from
$(1+x_{2r-3})y_{2r-3}$ and its square gives $1$ and $-2x_{2r-3}$.
Thus the two extractions reduce the filtration bound by two, and
\eqref{eq:dual-closed-path} shows that $g$ annihilates $P_{2r-1}^r$.
This gives a space of dimension $3^{r-2}$, supported on the interior
variables. It is independent of $\Omega_{2r-2}$, which annihilates the
entire open-path algebra. The dimensions already computed prove equality.
\end{proof}

\subsection{Odd cycles}
\label{sec:odd-cycles}

Use the pairing~\eqref{eq:coefficient-pairing} in the cyclic variables.
For $n=2m+1$ and $g\in F_{m+1}^\perp$, omit the two generators
$Y_{i-1},Y_i$ involving $\varphi_i$. The remaining generators form the
doubly-open path on the other $2m$ variables. Hence
Lemma~\ref{lem:dual-path} implies
\[
\partial_{\varphi_i}g\in\K\Omega_{2m-1}^{(i)},
\]
where the superscript means that $\varphi_i$ is omitted.
For all $i$, these conditions force every nonconstant monomial of $g$
to have degree $2m$, with all its coefficients equal. Thus
$g\in\K+\K\Omega_{2m}$.
Conversely, $\Omega_{2m}$ annihilates $\CC_{C_{2m+1}}$ by
\eqref{eq:open-path-algebra}. The constant polynomial $1$ annihilates
$F_{m+1}$: the identities
\begin{equation}\label{eq:cycle-adjacent-products}
Y_i^2=-2\varphi_i\varphi_{i+1},\qquad
Y_iY_{i+1}=\varphi_i\varphi_{i+1}-\varphi_i\varphi_{i+2}
                         +\varphi_{i+1}\varphi_{i+2}
\end{equation}
show that a repeated or adjacent pair has edge degree two. A product of
$m+1$ generators necessarily contains such a pair, because an independent
subset of $C_{2m+1}$ has at most $m$ vertices. Its edge degree is therefore
at most $2m$. Shorter products have the same bound. We have proved
\begin{equation}\label{eq:odd-cycle-annihilator}
F_{m+1}^\perp=\K\,1\oplus\K\Omega_{2m}.
\end{equation}
In other words, $F_{m+1}$ is the hyperplane in $\CC_{C_{2m+1}}$ whose
full-edge coefficient is zero. The product
\[
Y_0^2Y_1\prod_{j=1}^{m-1}Y_{2j+1}
\]
has length $m+2$ and a nonzero full-edge coefficient. This proves
Theorem~\ref{thm:cycle-top} for odd $n$.

\subsection{Even cycles and the closing kernel}
\label{sec:even-cycles}

Write $n=2m$, $a_j=\varphi_{2j}$, $b_j=\varphi_{2j+1}$
($j\in\Z/m\Z$). Identify square-free monomials in the $a_j$ with the
standard basis of the exterior algebra on $\K^m$, as vector spaces
only. Let $\iota$ be exterior contraction by the covector taking
value $1$ on every basis vector, and let
\[
\mathcal Z_{\rm odd}(a)=
\bigoplus_{\substack{k\ge1\\ k\ {\rm odd}}}
\ker\left(\iota:\bigwedge^k\K^m\longrightarrow
                         \bigwedge^{k-1}\K^m\right).
\]
Define $\mathcal Z_{\rm odd}(b)$ in the same way. Since the contraction
complex is exact, $\mathcal Z_{\rm odd}(a)$ is spanned by
\begin{equation}\label{eq:alternating-boundaries}
\iota a_J=
\sum_{\nu=1}^{2r}(-1)^{\nu-1}a_{J\setminus\{j_\nu\}},
\qquad J=\{j_1<\cdots<j_{2r}\},
\end{equation}
and
\begin{equation}\label{eq:odd-boundary-dimension}
\dim\mathcal Z_{\rm odd}(a)
=\sum_{k\ {\rm odd}}\binom{m-1}{k}=2^{m-2}.
\end{equation}

\begin{proposition}\label{prop:even-cycle-annihilator}
For $m\ge2$,
\begin{equation}\label{eq:even-cycle-annihilator}
F_m^\perp=
\K\Omega_{2m-1}\oplus\K(a_0+b_0)
\oplus\mathcal Z_{\rm odd}(a)\oplus\mathcal Z_{\rm odd}(b).
\end{equation}
\end{proposition}

For example, for $C_6$ ($m=3$) one has
$\mathcal Z_{\rm odd}(a)=\Span\{a_1-a_0,\,a_2-a_0\}$, and the
proposition says that $F_3^\perp$ is spanned by
\[
\Omega_5,\quad \varphi_0+\varphi_1,\quad
\varphi_2-\varphi_0,\quad\varphi_4-\varphi_0,\quad
\varphi_3-\varphi_1,\quad\varphi_5-\varphi_1 .
\]
Thus $\dim F_3=64-6=58$, and since $F_4=\CC_{C_6}$ has dimension $63$,
the top term of $H_6$ is $5t^4$, as in Theorem~\ref{thm:cycle-top}.
(A direct computation confirms these six annihilators
\cite[\texttt{cycles/annihilators.py}]{Cert}.)

We first restrict the support of a possible annihilator.
For $g\in F_m^\perp$, deleting $Y_{i-1},Y_i$ and applying
\eqref{eq:dual-odd-open-path} gives
\begin{equation}\label{eq:even-cycle-local-condition}
\partial_{\varphi_i}g\in
\K\Omega_{2m-2}^{(i)}+
(1-\varphi_{i+2})(1+\varphi_{i-2})\mathcal V_{m-1},
\end{equation}
with the second summand on the interior path
$\varphi_{i+2},\ldots,\varphi_{i-2}$.
This second summand lives on $2m-3$ variables, so its degree is at most
$2m-3$, while $\Omega^{(i)}_{2m-2}$ is homogeneous of degree $2m-2$.
Compare the two top degrees of $g$. If $g$ contains
$c\,\varphi_0\cdots\varphi_{2m-1}$, then $\partial_{\varphi_i}g$ has the
term $c\prod_{j\ne i}\varphi_j$ of degree $2m-1$, which does not occur on
the right of~\eqref{eq:even-cycle-local-condition}; hence $c=0$. Write
the degree-$(2m-1)$ part of $g$ as $\sum_kc_k\prod_{j\ne k}\varphi_j$.
Its extraction $\sum_{k\ne i}c_k\prod_{j\ne i,k}\varphi_j$ is the
degree-$(2m-2)$ part of $\partial_{\varphi_i}g$, so it is a multiple of
$\Omega^{(i)}_{2m-2}=\sum_{k\ne i}\prod_{j\ne i,k}\varphi_j$. Thus all
$c_k$ with $k\ne i$ are equal. Letting $i$ vary, all $c_k$ are equal,
and the degree-$(2m-1)$ part of $g$ is a multiple of $\Omega_{2m-1}$.
Subtract this multiple; it lies in $F_m^\perp$, since $\Omega_{2m-1}$
annihilates all of $\CC_{C_{2m}}$ by~\eqref{eq:open-path-algebra}. Condition
\eqref{eq:even-cycle-local-condition} now shows that no monomial of $g$
contains two cyclically adjacent variables.

Inductively suppose that no monomial contains two variables whose indices
differ by $1,3,\ldots,2s+1$. In the interior path of
\eqref{eq:even-cycle-local-condition}, use temporarily the local names
\[
 a_j=\varphi_{i+2+2j}\quad(0\le j\le m-2),\qquad
 b_j=\varphi_{i+3+2j}\quad(0\le j\le m-3).
\]
These names are local to this step, not the two global parity classes.
For $0\le s\le m-3$, extract the coefficient of
$b_s=\varphi_{i+2s+3}$.
In~\eqref{eq:dual-path-V}, $b_s$ occurs only as the middle alternative
in the $s$th factor. This extraction separates the remaining factors
into left and right groups. By induction the resulting coefficient is
independent of $a_0,\ldots,a_s,b_0,\ldots,b_{s-1}$. For $a_j$ the
index difference from $b_s$ is $2(s-j)+1$; for $b_j$ with $j<s$ the
index difference from $\varphi_i$ is $2j+3$. These are precisely among
the odd differences already excluded.
After setting the left $b$-variables to zero, its left dependence belongs
to
\begin{equation}\label{eq:no-constant-left-space}
(1-a_0)\Span\left\{
\prod_{j\in J}(a_j-a_{j+1}-a_ja_{j+1}):J\subseteq\{0,\ldots,s-1\}
\right\}.
\end{equation}
This space contains no nonzero constant. Indeed, every polynomial $P$
in the span has sum of linear coefficients zero, while
$(1-a_0)P=c$ would give $P=c(1+a_0)$ and hence $c=0$.
Applying this observation coefficientwise in the right variables excludes
the pair $\varphi_i,\varphi_{i+2s+3}$. Induction excludes all pairs of
opposite parity. For $m=2$, the initial adjacency restriction already
has this consequence. We have shown
\begin{equation}\label{eq:even-support-decomposition}
g=\lambda\Omega_{2m-1}+c+g_a(a)+g_b(b),
\end{equation}
where $g_a,g_b$ have zero constant term.

It remains to impose the closing condition on one parity class.
Choose one of $Y_{2j},Y_{2j+1}$ for each $j$, recording the choice by
$\varepsilon_j\in\{0,1\}$. Up to an overall sign, the coefficient of
$b_0\cdots b_{m-1}$ in the resulting product is
$\prod_j(1+u_ja_j)$, where $u_j=\varepsilon_{j-1}+\varepsilon_j-1$.
Define the linear map to functions on $\{0,1\}^m$ by
\begin{equation}\label{eq:Boolean-closing-map}
K_m(h)(\varepsilon)
=[a_0\cdots a_{m-1}]h(a)\prod_j(1+u_ja_j).
\end{equation}

\begin{lemma}\label{lem:Boolean-closing-kernel}
For $m\ge2$,
\begin{align}
\ker K_m&=\mathcal Z_{\rm odd}(a),\label{eq:Boolean-kernel}\\
\{h:K_m(h)\text{ vanishes on nonconstant }\varepsilon\}
 &=\mathcal Z_{\rm odd}(a)\oplus\K\,1\oplus\K a_0.
\label{eq:Boolean-two-exceptions}
\end{align}
\end{lemma}
\begin{proof}
Put $\sigma_j=2\varepsilon_j-1$, so $u_j=(\sigma_{j-1}+\sigma_j)/2$. The set $D$ of sign changes
of the cyclic sequence $\sigma$ is even, and every even subset occurs.
Also $u_j=0$ exactly on $D$. If $h=\sum_T h_Ta_T$, then
\[
K_m(h)(\sigma)=\sum_{T\supseteq D}h_T\prod_{j\notin T}u_j.
\]
Replacing $\sigma$ by $-\sigma$ separates the even- and odd-degree
parts of $h$. For the even part, take an even $D$ of maximal possible
degree with a nonzero coefficient. Its equation reads a nonzero multiple
of $h_D=0$. Descending induction shows that the even-degree part of
$h$ vanishes.
For the odd part, let $k$ be its highest degree and take $|D|=k-1$.
Only sets $T=D\cup\{j\}$ occur. Dividing by the common nonzero
product $\prod_{j\notin D}u_j$, the equation has coefficients $u_j$
for $j\notin D$. These signs change each time $j$ passes an element of
$D$, agreeing up to an overall sign with the coefficients of the
exterior contraction at $a_D$. Thus the equations are exactly
$\iota h_k=0$. For instance, if $D=\{d_1,d_2\}$, the sign is constant
on each complementary arc and opposite on the other arc, just as in
the coefficient of $a_{d_1}a_{d_2}$ in $\iota h_3$.

Every odd exterior-boundary cycle already lies in the kernel.
For $J=\{j_1<\cdots<j_{2r}\}$, direct expansion gives
\[
K_m(\iota a_J)
=\left(\prod_{j\notin J}u_j\right)
  \sum_{\nu=1}^{2r}(-1)^{\nu-1}u_{j_\nu}.
\]
If a sign change occurs outside $J$, the displayed product vanishes.
Otherwise $\sigma$ is
constant on each arc between successive elements of $J$, and the
alternating sum telescopes to zero. Thus the top odd part $h_k$ can be
subtracted, and descending induction proves~\eqref{eq:Boolean-kernel}.
Finally, $K_m(1)$ and $K_m(a_0)$ are supported only at the two constant
sign sequences and take independent values there. This proves
\eqref{eq:Boolean-two-exceptions}.
\end{proof}

\begin{proof}[Proof of Proposition~\ref{prop:even-cycle-annihilator}]
For a nonconstant choice sequence $\varepsilon$, the chosen generators
do not cover all the $a$-variables: a transition from $0$ to $1$ leaves
one $a_j$ absent. Hence their product pairs to zero with the constant
term and with $g_b$ in~\eqref{eq:even-support-decomposition}.
It follows that $K_m(g_a)$ vanishes at every nonconstant sequence;
the analogous argument applies to $g_b$.
Lemma~\ref{lem:Boolean-closing-kernel} now gives
\[
g=\lambda\Omega_{2m-1}+c+\mu a_0+\nu b_0+z_a+z_b,
\qquad z_a\in\mathcal Z_{\rm odd}(a),\quad
z_b\in\mathcal Z_{\rm odd}(b).
\]
Pairing with the two perfect-matching products
$\prod_jY_{2j}$ and $\prod_jY_{2j+1}$ gives
$c-\mu+\nu=0$ and $c+\mu-\nu=0$, respectively.
Thus $c=0$ and $\mu=\nu$.

Conversely, a polynomial in the $a$-variables can pair nontrivially with
a product of at most $m$ generators only if that product chooses exactly
one of $Y_{2j},Y_{2j+1}$ for every $j$. These tests are precisely
\eqref{eq:Boolean-closing-map}; therefore every $z_a$ and $z_b$
annihilates $F_m$. To test $a_0+b_0$, products of fewer than $m$
generators have edge degree at most $2m-2$. For $m$ factors the same is
true if there is a repeated or adjacent pair, by
\eqref{eq:cycle-adjacent-products}. The only remaining products are the
two perfect matchings, and their tests give zero. Finally
$\Omega_{2m-1}$ annihilates the whole algebra.
The summands are independent: the linear parts of the two $\mathcal Z$
spaces have coefficient sum zero separately on each parity class, unlike
$a_0+b_0$; and $\Omega_{2m-1}$ has mixed support and higher degree.
\end{proof}

\begin{proof}[Completion of the proof of Theorem~\ref{thm:cycle-top}]
For $n=2m$, Proposition~\ref{prop:even-cycle-annihilator} and
\eqref{eq:odd-boundary-dimension} give
$\dim F_m^\perp=2^{m-1}+2$.
At filtration degree $m+1$, the coefficient extractions instead
annihilate the exhausted open paths $P_{2m-1}^{m+1}$.
Thus $F_{m+1}^\perp\subseteq\K+\K\Omega_{2m-1}$.
A perfect-matching product already has nonzero full-edge coefficient,
so the constant summand is excluded. Hence
$F_{m+1}^\perp=\K\Omega_{2m-1}$ and
$F_{m+1}=\CC_{C_{2m}}$. The top coefficient is
$(2^{m-1}+2)-1=2^{m-1}+1$.
\end{proof}

\subsection{The cycle conjecture}
\label{sec:cycle-conjecture}

Theorem~\ref{thm:cycle-top} fixes the top term of $H_n$, and
Theorem~\ref{thm:cycle-low-coefficients} proves the first three stable
coefficients. The conjecture extends this agreement to all degrees below
half the cycle length, using the trace of the transfer matrix
$M=M(t)$ of~\eqref{eq:propagation-matrix}. Put
$\mathrm{tr}_n(t)=\operatorname{tr}M(t)^n$.

The trace has the following weighted independent-set interpretation. For a set $S$ of vertices
of $C_n$, let $c(S)$ be the number of vertices outside $S$ both of whose
neighbors lie in $S$.

\begin{proposition}[A weighted independence form]
\label{prop:cycle-weighted-independence}
For $w\in\K$, let $M_w(t)$ be the matrix~\eqref{eq:propagation-matrix}
with the entry $2$ replaced by $w-1$, so that $M_3=M$. For $n\ge3$,
\[
 \operatorname{tr}M_w(t)^n=\sum_{S}w^{c(S)}t^{|S|},
\]
the sum over the independent sets $S$ of vertices of $C_n$. In
particular $\mathrm{tr}_n(t)=\sum_S3^{c(S)}t^{|S|}$, and
$\operatorname{tr}M_1(t)^n=I(C_n;t)$.
\end{proposition}
\begin{proof}
The potentially nonzero transition weights of $M_w$ are $w-1$ from $e_0$ to $e_1$, $t$
from $e_1$ to $e_0$, $1$ from $e_1$ to $e_2$, $t$ from $e_2$ to $e_1$
and $1$ from $e_2$ to $e_2$. Thus $\operatorname{tr}M_w^n$ is the sum,
over the sequences $(s_i)_{i\in\Z/n\Z}$ of states, of the products of the
entries from $s_i$ to $s_{i+1}$. Put $S=\{i:s_i=e_1\}$ and
$D=\{i:s_i=e_0\}$. There is no entry from $e_1$ to $e_1$, so $S$ is
independent, and $e_0$ is entered and left only through $e_1$, so every
element of $D$ has both neighbors in $S$. Conversely, every pair $(S,D)$
of this kind comes from exactly one sequence, with $s_i=e_2$ outside
$S\cup D$. Each element of $D$ contributes the factor $t$ of its
entering step; each element of $S$ contributes $w-1$ if it is entered
from $D$ and $t$ otherwise. Since the elements of $S$ entered from $D$
correspond to the elements of $D$, the product is
$t^{|S|}(w-1)^{|D|}$, and summing over $D$ gives $t^{|S|}w^{c(S)}$.
\end{proof}

\begin{conjecture}[Cycle congruence]
\label{conj:cycles}
Let $\alpha\ne0$ and $H_n=\hs^f_{C_n}$. For every $n\ge3$,
\begin{equation}\label{eq:cycle-stable-congruence}
H_n(t)\equiv \mathrm{tr}_n(t)-t^2\pmod{t^{\lceil n/2\rceil}}.
\end{equation}
Equivalently, for every $d\ge0$ with $2d<n$,
\[
 [t^d]H_n=\sum_{\substack{S\subseteq V(C_n)\text{ independent}\\|S|=d}}
            3^{c(S)}-\delta_{d,2}.
\]
Thus each independent set is weighted by a factor of three for each
vertex outside it whose two neighbors lie in the set.
\end{conjecture}

For $d=2$ and $n\ge5$, the weighted sum is
$3n+n(n-5)/2=\binom{n+1}{2}$, consistent with the stable quadratic
coefficient. Setting $w=1$ gives $I(C_n;t)$, which differs by $t$ from
the matching-quotient series of Remark~\ref{rem:cycle-matching-quotient}.
At $t=1$ the weighted sum is $2^n+(-1)^n\ell_n$, as used below.

The conjecture agrees with the series $H_n$ for $n\le17$, computed
modulo two primes \cite[\texttt{cycles/}]{Cert} by the method of
Proposition~\ref{prop:cycle-interpolation}.
It was first found in the following equivalent form, as a recurrence for
the operator $L$ of~\eqref{eq:L}. Here $\ell_k$ are the Lucas numbers
defined in Proposition~\ref{prop:degenerate-cycles}.

\begin{proposition}[Equivalent forms of the conjecture]
\label{prop:cycle-reduction}
Given Theorem~\ref{thm:cycle-top}, Conjecture~\ref{conj:cycles} is
equivalent to the recurrence
\begin{equation}
    \label{eq:circ-rec}
    L(H_n)=\begin{cases}
      t^3+ (2^{n/2-2}+1)t^{n/2}(1-t) & \text{if $n$ is even},   \\
      t^3+ (\ell_{n-4}-1+2t)(1-t)t^{(n-1)/2} & \text{if $n$ is odd,}
    \end{cases}
    \qquad n\ge6,
\end{equation}
together with the initial values
\[
H_3=1+3t+2t^2+t^3, \quad
H_4=1+4t+7t^2+3t^3,  \quad
H_5=1+5t+14t^2+10t^3+t^4.
\]
It is also equivalent to the closed formulas~\eqref{eq:cycle-even-trace}
and~\eqref{eq:cycle-odd-trace} below.
\end{proposition}
\begin{proof}
The characteristic polynomial of $M$ is
$z^3-z^2-3tz+2t$ by~\eqref{eq:propagation-matrix}, so Cayley--Hamilton
gives $L(\mathrm{tr}_n)=0$, that is, $\mathrm{tr}_n=\mathrm{tr}_{n-1}+3t\mathrm{tr}_{n-2}-2t\mathrm{tr}_{n-3}$ for $n\ge3$.
The initial values $\mathrm{tr}_0=3$, $\mathrm{tr}_1=1$ and
$\mathrm{tr}_2=1+6t$, together with this recurrence, give
$\deg\mathrm{tr}_n\le\lfloor n/2\rfloor$ by induction. Moreover
\[
\det(zI-M(1))=(z-2)(z^2+z-1),\qquad
\mathrm{tr}_n(1)=2^n+(-1)^n\ell_n.
\]
Since $\dim\CC_{C_n}=2^n-1$ by~\eqref{eq:open-path-algebra},
the congruence and the top-term theorem leave just one coefficient
undetermined, and evaluation at $t=1$ determines it.
Writing $A_m=2^{m-1}+1$, the resulting formulas are
\begin{align}
H_{2m}&=\mathrm{tr}_{2m}-t^2-(\ell_{2m}+A_m)t^m+A_mt^{m+1}
 &&(m\ge2),\label{eq:cycle-even-trace}\\
H_{2m+1}&=\mathrm{tr}_{2m+1}-t^2+(\ell_{2m+1}-1)t^{m+1}+t^{m+2}
 &&(m\ge1).\label{eq:cycle-odd-trace}
\end{align}
Substitution in $L$, using $\ell_n=\ell_{n-1}+\ell_{n-2}$,
gives the two lines of the conjectured recurrence. The formulas also give
its initial values. Conversely, that recurrence and those initial values
uniquely determine the sequence, so they imply
\eqref{eq:cycle-even-trace}--\eqref{eq:cycle-odd-trace}, and hence the
congruence.
\end{proof}

\begin{remark}[The matching quotient and the Lucas contribution]\label{rem:cycle-matching-quotient}
The number $\ell_n$ counts all independent vertex subsets of $C_n$.
There is an unconditional occurrence of this count: if
$J=(\varphi_i\varphi_{i+1}:i\in\Z/n\Z)\subset\Phi_n$, then
$J\subset\CC_{C_n}$ and the induced quotient filtration has series
$I(C_n;t)-t$, where $I(C_n;t)$ is the independence polynomial.
Indeed, modulo $J$ the generators become linear differences, and
\eqref{eq:open-path-algebra} contains every surviving monomial of degree
at least two and all but one linear direction. Thus this quotient has
dimension $\ell_n-1$ for both parities.

This quotient is the Dunkl matching algebra of the cycle in
Grinberg--Kirillov--Shapiro \cite{GKS}. Their matching algebra
$\mathcal M(G)$ is the quotient of the square-zero edge algebra by
products of incident edges; $\mathcal D(G)$ is its subalgebra generated
by the signed incidence sums. In $\mathcal M(G)$ every $X_v^2=0$, so
all normalized generators $f(X_v)$ have the same image $X_v$.
For $G=C_n$, the additional ideal is precisely $J$, and hence
$\CC_{C_n}^f/J\simeq\mathcal D(C_n)$ with its degree filtration.
Their unicyclic formula \cite[Theorem~5.16]{GKS} gives the same series
$I(C_n;t)-t$ (matchings of a cycle correspond to independent sets of
its line graph, again a cycle). This identifies the quotient, not the
induced filtration on $J$, and therefore does not prove the full
cycle formula.

In the conjectural trace formulas the Lucas part is
$b_n=(-1)^{n+1}\ell_nt^{\lceil n/2\rceil}$. Directly from the Lucas
recurrence,
\[
L(b_n)=0\ (n\text{ even}),\qquad
L(b_n)=\ell_{n-4}t^{(n-1)/2}(1-t)\ (n\text{ odd}).
\]
Thus the even recurrence has no explicit Lucas term because of a
cancellation, although the independent sets contribute for both
parities.
\end{remark}

\begin{remark}[An exact boundary-operator formulation]\label{rem:cycle-boundary-operator}
For $2d<n$, every product of at most $d$ cyclic generators omits two
consecutive generators. Let $P_{n-1}^{(i),d}$ be the open path obtained
by omitting $Y_{i-1},Y_i$, and put
$E_{i,d}=(P_{n-1}^{(i),d})^\perp$ in the algebra on the other variables.
Then
\[
F_d^\perp=\ker D_{n,d},\qquad
D_{n,d}:\Phi_n\longrightarrow
\bigoplus_i\K[\varphi_j:j\ne i]/(\varphi_j^2)\big/E_{i,d},
\quad g\longmapsto([\partial_{\varphi_i}g])_i.
\]
Consequently~\eqref{eq:cycle-stable-congruence} is equivalent to
\[
\rank D_{n,d}=\sum_{j=0}^d[t^j]\mathrm{tr}_n-\mathbf1_{\{d\ge2\}},
\qquad 2d<n.
\]
This rank identity is open. It encodes the condition that the
separate path coefficients come from one polynomial, which the trace of
the transfer matrix alone does not see.
\end{remark}


\begin{proposition}[An interpolation matrix for the cycle filtration]
\label{prop:cycle-interpolation}
For every $n\ge3$ and $d\ge0$, let $A_{n,d}$ be the integer matrix with
rows indexed by $S\subseteq\Z/n\Z$ and columns indexed by
$m\in\Z_{\ge0}^n$ with $|m|\le d$, and entries
\begin{equation}\label{eq:cycle-interpolation-matrix}
 (A_{n,d})_{S,m}=\prod_{i\in S}(m_{i-1}-m_i).
\end{equation}
Then $\dim F_d=\rank_{\mathbb Q}A_{n,d}$.
\end{proposition}
\begin{proof}
In the normalization~\eqref{eq:cycle-normalization},
$Y_i=\exp(\varphi_{i+1}-\varphi_i)-1$ because the cube of that difference
is zero. Replacing the generators $Y_i$ by $1+Y_i$ preserves every
word-filtration level: binomial expansion gives inclusions in both
directions without increasing the number of factors. Thus $F_d$ is
spanned by
\[
 \prod_i(1+Y_i)^{m_i} =\exp\!\left(\sum_i(m_{i-1}-m_i)\varphi_i\right)
 =\prod_i\bigl(1+(m_{i-1}-m_i)\varphi_i\bigr),\qquad |m|\le d.
\]
Their coefficients in the square-free basis are exactly~\eqref{eq:cycle-interpolation-matrix}.
Since these coefficients are integers, their rank is the same over every characteristic-zero field.
\end{proof}

For $2d<n$, Conjecture~\ref{conj:cycles} is therefore equivalent to
\[
 \rank_{\mathbb Q}A_{n,d}
 =\sum_{j=0}^d[t^j]\mathrm{tr}_n-\mathbf1_{\{d\ge2\}}.
\]
The unresolved step is the rank of these columns, which retain the
joint cyclic compatibility. Theorem~\ref{thm:cycle-low-coefficients}
settles it through $d=3$ in the stable range.

